% \begin{document}
\chapter{Teoria dei Caratteri}
\section{Caratteri}
Abbiamo ora uno strumento potentissimo per \textbf{classificare} le rappresentazioni senza lavorare esplicitamente con le matrici: i \textbf{caratteri}.

\dfn{Traccia}{
Sia $V$ uno spazio vettoriale di dimensione finita e $f: V \to V$ un'applicazione lineare. In qualunque base di $V$, la matrice di $f$ è $A = (a_{ij})$. La traccia di $f$ è:

$$\mathrm{tr}(f) := \sum_i a_{ii}$$

cioè la somma degli elementi diagonali.
}

Le \textbf{proprieta' fondamentali} della traccia sono:
\begin{enumerate}
\item $\mathrm{tr}(\mathrm{Id}_V) = \dim V$ — è uno dei coefficienti del polinomio caratteristico.
\item \textbf{Non dipende dalla base}: se $A$ e $B = P^{-1}AP$ sono la stessa mappa in basi diverse, allora $\mathrm{tr}(B) = \mathrm{tr}(P^{-1}AP) = \mathrm{tr}(A)$ (per ogni $A$ invertibile, $\mathrm{tr}(PQ) = \mathrm{tr}(QP)$).
\item $\mathrm{tr}(f) = \sum_i \lambda_i$ dove $\lambda_i$ sono gli autovalori di $f$ contati con molteplicità (radici del polinomio caratteristico).
\end{enumerate}
La proprietà 2 è cruciale: la traccia è un \textbf{invariante} dell'operatore, indipendente dalla scelta della base.

\dfn{Carattere}{
  Sia $(V, \rho)$ una rappresentazione di $G$ su $\mathbb{C}$. Il carattere di $\rho$ è la funzione:

$$\chi_\rho : G \longrightarrow \mathbb{C}, \qquad g \longmapsto \mathrm{tr}(\rho(g))$$
}
\subsection{Proprieta' dei Caratteri}
Le proprieta' piu' interessanti dei caratteri sono:
\begin{enumerate}
\item $\chi_\rho(e) = \dim V$.
\item $\chi_\rho$ è costante sulle classi di coniugio di $G$:
  $$g' = hgh^{-1} \implies \chi_\rho(g') = \chi_\rho(g).$$
    Una funzione costante sulle classi di coniugio si chiama funzione di classe. I caratteri sono dunque funzioni di classe — questo è fondamentale: non dobbiamo calcolare $\chi$ su tutti gli elementi del gruppo, basta una volta per ogni classe di coniugio.
\item $\chi_\rho(g^{-1}) = \overline{\chi_\rho(g)}$ (il complesso coniugato), per $G$ finito.
\end{enumerate}

\nt{
  Gli autovalori di $\rho(g)$ sono radici dell'unità (perché $g^n = e$ per qualche $n$, quindi $\rho(g)^n = I$), quindi hanno modulo 1, e gli autovalori di $\rho(g^{-1}) = \rho(g)^{-1}$ sono i loro inversi, che coincidono con i coniugati.
}

\pf{Dimostrazione delle proprieta'}{
  Dimostriamo in ordine:
  \begin{enumerate}
  \item $\rho(e) = \mathrm{Id}_V$, e $\mathrm{tr}(\mathrm{Id}_V) = \dim V$.
    \item $\chi_\rho(hgh^{-1}) = \mathrm{tr}(\rho(h)\rho(g)\rho(h)^{-1}) = \mathrm{tr}(\rho(g)) = \chi_\rho(g)$, usando ciclicità della traccia.
  \end{enumerate}
}

\ex{Azione naturale di $ S_3 $ su $ \mathbb{C}^3 $}{
  Sia $G = S_3$ con la rappresentazione naturale su $V = \mathbb{C}^3$ con base ${e_1, e_2, e_3}$:

$$g \cdot (x_1, x_2, x_3) = (x_{g(1)}, x_{g(2)}, x_{g(3)})$$

Le classi di coniugio di $S_3$ sono tre: ${e}$, le trasposizioni ${(12),(13),(23)}$, i 3-cicli ${(123),(132)}$.\\

Calcolo di $\chi$ per ciascuna classe:
\begin{itemize}
\item $g = \mathrm{id}$: $\rho(e) = I_3$, quindi: $$\chi(e) = \mathrm{tr}(I_3) = 3$$
  \item $g = (12)$: la matrice permuta le prime due coordinate, fissa la terza: $$\rho((12)) = \begin{pmatrix}0&1&0\\1&0&0\\0&0&1\end{pmatrix}, \qquad \chi((12)) = 0+0+1 = 1$$

Lo stesso vale per tutte le trasposizioni: ciascuna fissa esattamente un indice, quindi la traccia è sempre 1.
\item $g = (123)$: la ciclica $1\mapsto 2\mapsto 3\mapsto 1$, nessun punto fisso: $$\rho((123)) = \begin{pmatrix}0&0&1\\1&0&0\\0&1&0\end{pmatrix}, \qquad \chi((123)) = 0+0+0 = 0$$

Lo stesso per $(132)$: nessun punto fisso, traccia 0.
\end{itemize}
\textbf{Regola generale}: per la rappresentazione di permutazione, $\chi(g) = $ numero di punti fissi di $g$.

TODO: fai tavola

I valori sono interi — non è un caso: vedremo che i caratteri di gruppi simmetrici sono sempre interi. I valori $3, 1, 0$ codificano interamente la struttura della rappresentazione. Nella prossima lezione vedremo come usarli per decomporre $V$ in irriducibili tramite le relazioni di ortogonalità tra caratteri.
}
\section{Richiami di algebra lineare}

Dati due spazi vettoriali $V$ e $U$ su $K$, conosciamo già come costruire la somma diretta $V \oplus U$. Vogliamo ora costruire un altro oggetto fondamentale: il prodotto tensoriale $V \otimes U$.

Ripassino per informatici di somma diretta:
\dfn{Somma diretta}{
  Dati due spazi vettoriali $V$ e $U$ su $K$, la somma diretta $V \oplus U$ tale spazio vettoriale:
  \[
  V \oplus U := V \times U \text{ con operazioni per coordinate}
  \]
  Sia ${v_i}_{i\in I}$ e ${u_j}_{j\in J}$ due base di $V$ e $U$ rispettivamente, allora
  \[
  \{(v_i,0)\}_{i\in I}\cup\{(0,u_j)\}_{j\in J} \text{ è una base di } V \oplus U.
  \]
}

\cor{}{
  $V \oplus U$ ha dimensione $\dim V + \dim U$.
}

\subsection{Prodotto tensoriale}
Nella somma diretta le due componenti "vivono separatamente". Il prodotto tensoriale serve invece a modellare le applicazioni bilineari: oggetti che dipendono linearmente da ciascuno dei due argomenti separatamente, ma non sono lineari nella coppia

Si vuole infatti definire uno spazio $V \otimes U$ tale che dare un'applicazione lineare $V \otimes U \to W$ è equivalente a dare un'applicazione bilineare $V \times U \to W$

\subsubsection{Costruzione}

TODO: non penso sia giusta questa definizione. Gem sembra confondersi perche' preferisce usare $ F(V \otimes U) $ per indicare lo SV libero, mentre usa $ \mathcal{L}(V,U) $ per indicare gli omomorfismi fra $ V $ e $ U $, ma non e' questo che vogliamo. Comunque lo tengo perche' forse non ho capito io. 

Va bene buon Basta, cambio la definzione, specificatamente per $U\times V$ (ti amo)

Qui definisco ciò, ma è sonnet la fonte sta volta, se tu hai fonte diversa dimmi pure:

\dfn{Spazio vettoriale libero}{
  Sia $S$ un insieme qualsiasi e $K$ un campo. Lo spazio vettoriale libero su $S$ su $K$ è:

  \[
    L(S):=\left\{f:S\to K\;\middle|\;f(s)=0\text{ per quasi ogni } s\in S\right\}
  \]
  dove "quasi ogni" significa "per tutti tranne un numero finito".
}

\dfn{Impulso di Dirac}{
  Per ogni $s \in S$ \textit{Impulso di Dirac} la seguente funzione:
  \[
    \delta_s : S \to \mathbb{N}, \qquad \delta_s(t) = \begin{cases} 1 & \text{se } t = s \\ 0 & \text{se } t \neq s \end{cases}
  \]
}

Si può dimostrare che questa sarà la base di ciò che noi chiameremo spazio vettoriale libero di $U\times V$

\dfn{Spazio vettoriale libero di $U\times V$}{
  Siano $V$ e $U$ spazi vettoriali su $K$. Prendiamo $S=V\times U$, l'insieme di tutte le coppie $(v,u)$ con $v\in V$ e $u\in U$.

  Allora 
  $$L(V\times U) = \left\{ \sum_{\text{finita}} \alpha_{(v,u)} \cdot \delta_{(v,u)} \;\middle|\; \alpha_{(v,u)} \in K \right\}$$
  La base è $\{\delta_{(v,u)}\}_{(v,u) \in V \times U}$, una per ogni coppia.
}

\nt{
  Si noti che $(v,u)$ è una coppia che rappresenta tutte le possibili coppie tra $V$ e $U$, non le base. PROPRIO LE COPPIE CAZZO 
}

\nt{Si noti che l'impulso di Dirac è una base di $L(V\times U)$ come definito prima}

\ex{Esempietto}{
  Sia $V=U=\mathbb{R}^2$ con basi $\{e_1,e_2\}$ e $\{f_1,f_2\}$.

Alcuni elementi di base di $L(V\times U)$ sono:

$$\delta_{(e_1, f_1)}, \quad \delta_{(e_2, f_2)}, \quad \delta_{(e_1 + e_2,\; f_1)}, \quad \delta_{(3e_1,\; 2f_2)}, \quad \delta_{(0,\; f_1)}, \quad \ldots$$

Un elemento generico è:
$$
w=5\,\delta_{(e_1, f_1)} - 2\,\delta_{(e_1+e_2,\; f_2)} + 7\,\delta_{(3e_1,\; f_1)}
$$
}

\nt{
  In $L(U\times V)$non vale nessuna relazione algebrica tra le coppie. In particolare queste tre cose sono veri e propri assiomi di $L(V\times U)$:

  $$\delta_{(e_1 + e_2,\; f_1)} \neq \delta_{(e_1, f_1)} + \delta_{(e_2, f_1)}$$

  $$\delta_{(2e_1,\; f_1)} \neq 2\,\delta_{(e_1, f_1)}$$

  $$\delta_{(0,\; f_1)} \neq 0$$
}

\cor{}{
  Ogni elemento $\alpha \in L(V, U)$ può essere scritto in modo unico come una combinazione lineare formale finita:
  $$\alpha = \sum_{i=1}^n c_i (v_i, u_i)$$
  dove i coefficienti $c_i$ sono scalari e le coppie $(v_i, u_i)$ sono gli elementi della base
}

Adesso definiamo un sottospazio
\dfn{}{
  Sia $B \subseteq L(V, U)$ il sottospazio generato da tutte le relazioni di bilinearità, cioè da tutti gli elementi della forma:

  $$
  \begin{aligned}
    &(a_1 v_1 + a_2 v_2,, u) - (a_1(v_1, u) + a_2(v_2, u)) \\
    &(v,, b_1 u_1 + b_2 u_2) - (b_1(v, u_1) + b_2(v, u_2))
  \end{aligned}
  $$
per ogni $v,v_1,v_2 \in V$ e $u,u_1,u_2 \in U$ e $a_1,a_2,b_1,b_2 \in K$.
}
\nt{
  Il sottospazio $ B $ e' quindi la combinazione lineare di tutti i vettori ottenuti prendendo per ogni vettore in $ V $ e in $ U $ la forma lineare (a sx e a dx) e sottraendo il vettore che, \textbf{se valesse la bilinearita'}, sarebbe uguale.\\

  Quindi stiamo essenzialmente creando lo spazio vettoriale contenente tutti i vettori che noi vorremmo fossero nulli.
}

Possiamo quindi definire il prosotto tensoriale come:

\dfn{Prodotto tensoriale}{
  Il prodotto tensoriale di $V$ e $U$ è:

  $$V \otimes U := L(V, U) / B$$
}
Si noti questo teorema bellissimo degli insiemi quozienti

\thm{Struttura Vettoriale del Quoziente}{
  Se $V$ è uno spazio vettoriale e $W$ è un suo sottospazio, la relazione $\xi \sim \zeta \iff \xi - \zeta \in W$ definisce uno spazio vettoriale quoziente $V/W$.
  Le operazioni di somma e prodotto per scalare sono ben definite dalle relazioni:
  \[
  [v] + [u] = [v + u]
  \]
  \[
  \lambda [v] = [\lambda v]
  \]
}
\pf{
  Sketch della dimostrazione
}{
    La dimostrazione della "buona definizione" consiste nel verificare che, se scegliamo rappresentanti diversi per la stessa classe (ovvero $v \sim v'$ e $u \sim u'$), il risultato della somma non cambia. Poiché $(v' + u') - (v + u) = (v' - v) + (u' - u)$, e sapendo che entrambi i termini a destra appartengono a $W$, la loro somma appartiene ancora a $W$. Pertanto $[v' + u'] = [v + u]$.
}

Andiamo inoltre a mostrare altre cose dell'insieme quoziente
\dfn{0 quoziente}{
In uno spazio quoziente $L/B$, gli elementi sono "classi di equivalenza" (insiemi di vettori). La classe dello zero è definita come:

$$[0]=\{w\in L\mid w-0\in B\}=\{w\in L\mid w\in B\}=B$$


Quindi, lo zero del quoziente è il sottospazio $B$ stesso.
}

\nt{
  Qualsiasi vettore che in $L(V\times U)$ apparteneva a $B$, una volta passato al quoziente, diventa indistinguibile dallo zero. Ovvero:
  \[
  \delta_{(v+v', u)} - \delta_{(v,u)} - \delta_{(v',u)} \in B \implies [\delta_{(v+v', u)}] = [\delta_{(v,u)}] + [\delta_{(v',u)}]
  \]
}

Si possono così definire i cosiddetti tensori semplici
\dfn{Tensori semplici}{
  Si definisce \textit{Tensore semplice} il seguente elemento:
  \[
    v \otimes u := [(v,u)]
  \]
}



\ex{Polinomi in due variabili}{
$V = K[x]$, $U = K[y]$ (polinomi in $x$ e in $y$ rispettivamente). Le basi sono $\{1, x, x^2, \ldots\}$ e $\{1, y, y^2, \ldots\}$. Il prodotto tensoriale ha base $\{x^i \otimes y^j\}_{i,j \in \mathbb{N}}$, cioè:

$$K[x] \otimes_K K[y] \cong K[x, y]$$

l'anello dei polinomi in due variabili.
}

\ex{Algebra del prodotto diretto}{
  Siano $G, H$ gruppi finiti. Allora:

$$K[G \times H] \cong K[G] \otimes_K K[H]$$

L'algebra di gruppo del prodotto diretto è il prodotto tensoriale delle singole algebre di gruppo.
}

\subsection{Spazio Hom e Spazio Duale}

\dfn{Sapzio di Hom}{
  Dati $V, U$ spazi vettoriali su $K$, lo spazio delle applicazioni lineari da $V$ a $U$ è:

  \begin{equation*}
    \mathrm{Hom}_K(V, U) := \{f : V \to U \mid f \text{ lineare}\}
  \end{equation*}

  con le ovvie operazioni puntali. $\mathrm{Hom}_K(V, U)$ è uno spazio vettoriale su $K$ di dimensione $\dim V \cdot \dim U$ (corrisponde alle matrici $(\dim U) \times (\dim V)$).
}

Dai cui discende 
\dfn{Spazio Duale}{
  Lo spazio duale di $V$ è:
  $$V^* := \mathrm{Hom}_K(V, K)$$
  cioè lo spazio delle forme lineari su $V$ (funzionali lineari a valori in $K$). Si ha $\dim V^* = \dim V$
}


\dfn{Isomorfismo canonico}{
  Sia $V, U$ spazi vettoriali su $K$. Definiamo:
  \[
    \Psi:V^* \otimes U \to \mathrm{Hom}_K(V, U), \quad \varphi \otimes u \mapsto (v \mapsto \varphi(v) \cdot u)
  \]
}
\pf{Dim}{
  Lasciata al basta\\

  Ma e' una definizione :cryingemoji: cosa devo dimostrare?
}

\mprop{
  Ben definizione dell'isomorfismo canonico
}{
  La funzione definita in precedenza e' ben definita quindi:
  \[
    V^* \otimes U \cong \mathrm{Hom}_K(V, U)
  \]
  (se $ V $ e $ U $ sono finiti)
}

\pf{Dimostrazione}{
  Dobbiamo dimostrare la bilinearita' di $ \varphi(v)\cdot u $ su $ (\varphi, u) $. Usando le definizioni di somma e prodotto nello spazio duale e le proprieta' di linearita' e del prodotto scalare in $ U $ si puo' dimostrare che
  \[
  [(a\varphi_1+b\varphi_2)(v)]\cdot u = a(\varphi_1(v)\cdot u)+b(\varphi_2(v)\cdot u)
  \]
  e 
  \[
  \varphi(v)\cdot(cu_1+du_2) = c(\varphi(v)\cdot u_1)+d(\varphi(v)\cdot u_2)
  \]
  Per la \textbf{proprieta' universale del prodotto tensoriale}, la bilinearita' della definizione di $ \Psi $ su $ V^* \times U $ implica che $ \Psi $ e' un'applicazione lineare (quindi ben definita) da $ V^* \otimes U $, e quindi che rispetta le classi di equivalenza.\\

  Ora dimostriamo l'isomorfismo. Dato che $ dim(V^* \otimes U) = n \cdot m = dim(Hom(V,U)) $, ci basta dimostrare l'iniettivita' per ottenere anche la suriettivita' (o viceversa).

  Dimostrare l'iniettivita' di un'applicazione lineare equivale a dimostrare che
  \[
    ker(\Psi) = \{n\}
  \]
  dove $ n $ e' il tensore nullo. Questo si puo' fare dimostrando che se 
  \[
    \Psi(t) = N
  \]
  dove $ N $ e' l'operatore nullo, allora applicando i vettori base $ e_1,...,e_n \in V $ dobbiamo ottenere il vettore nullo in $ U $ (dato che l'operatore e' un omomorfismo):
  \[
    \Psi(t)(e_k) = 0
  \]
  Si puo' dimostrare che questo e' vero solo se i coefficenti di $ t $ sono tutti zero, ovvero $ t = n $.
}

\section{Costruzione di Rappresentazioni complesse e i loro Caratteri}

Si può dimostrare che Se $V$ e $U$ sono rappresentazioni di $G$, anche $V \oplus U$, $V \otimes U$, $V^*$ e $\mathrm{Hom}_K(V, U)$ sono naturalmente rappresentazioni di $G$.


Siano $\rho:G\to GL(V)$ e $\pi:G\to GL(U)$ due rappresentazioni
\subsection{Somma diretta di rappresentazioni}
Vogliamo dare a $ V \oplus U $ una struttura di rappresentazione, definiamo quindi $ "\rho \oplus \pi": V \oplus U \to V \oplus U $ come:
\[
  (\rho \oplus \pi)(g)(v + u) = \rho(g)(v) + \pi(g)(u)
\]

Si verifica subito che è un omomorfismo.

In termini matriciali: scelta la base $\{v_1, \ldots, v_n\}$ di $V$ e $\{u_1, \ldots, u_m\}$ di $U$, la matrice di $(\rho \oplus \pi)(g)$ nella base unione è \textbf{a blocchi diagonali}:

$$(\rho \oplus \pi)(g) \longleftrightarrow \begin{pmatrix} A & 0 \\ 0 & B \end{pmatrix}$$

dove $A = \rho(g)$ e $B = \pi(g)$ nelle rispettive basi.

\mprop{Carattere di una somma diretta}{
  Date due rappresentazioni $ V, U $ con caratteri $ \chi_V, \chi_U $:
  \[
    \chi_{V \oplus U}(g) = \chi_V(g) + \chi_U(g) 
  \]
}
\pf{Dimostrazione}{
  Ovvia data la matrice associata.
}

\subsection{Prodotto tensoriale di rappresentazioni}
In modo analogo, definiamo $ "\rho \otimes \pi": G \to GL(V \otimes U) $ come:

$$(\rho \otimes \pi)(g)(v \otimes u) = (\rho(g)v) \otimes (\pi(g)u)$$

Nella base $\{v_i \otimes u_j\}$, la matrice di $(\rho \otimes \pi)(g)$ è il \textbf{prodotto di Kronecker} delle matrici $\rho(g)$ e $\pi(g)$.

\mprop{Carattere di un prodotto tensoriale}{
  Date due rappresentazioni $ V, U $ con caratteri $ \chi_V, \chi_U $:
  \[
    \chi_{V \otimes U}(g) = \chi_V(g) \cdot \chi_U(g) 
  \]
}
\pf{Dimostrazione}{
  Nella base $\{v_i \otimes u_j\}$, gli elementi diagonali della matrice di $(\rho \otimes \pi)(g)$ corrispondono agli indici $(i,j) = (i,j)$. Esplicitamente:

$$\mathrm{tr}(\rho \otimes \pi)(g) = \sum_{i=1}^n \sum_{j=1}^m [\rho(g)]_{ii} \cdot [\pi(g)]_{jj} = \left(\sum_i [\rho(g)]_{ii}\right)\left(\sum_j [\pi(g)]_{jj}\right) = \mathrm{tr}(A)\cdot\mathrm{tr}(B).$$

}

\subsection{Spazio duale di una rappresentazione}
Vogliamo una rappresentazione $\rho^*: G \to \mathrm{GL}(V^*)$ che sia compatibile con la struttura. La condizione naturale è la \textbf{$G$-invarianza} del pairing: vogliamo che

$$\langle \varphi, v \rangle = \langle g \cdot \varphi,\; g \cdot v \rangle = \langle \rho^*(g)(\varphi),\; \rho(g)(v) \rangle$$

cioè il "prodotto scalare naturale" $\langle \varphi, v \rangle := \varphi(v)$ sia invariante. Questo forza:

$$\rho^*(g)(\varphi)(v) = \varphi(\rho(g)^{-1} v)$$

In altre parole: $g$ agisce su $V^*$ tramite \textbf{l'inverso della trasposta} di $\rho(g)$, cioè $\rho^*(g) = (\rho(g)^{-1})^T = (\rho(g^{-1}))^T$.

\mprop{Carattere del duale}{
  Data una rappresentazione $ V$ con carattere $ \chi_V$:
  \[
    \chi_{V^*}(g) = \overline{\chi_V(g)}
  \]
}
\pf{Dimostrazione}{
  Gli autovalori di $\rho(g)$ sono radici dell'unità $\lambda_1, \ldots, \lambda_n$ (perché $g$ ha ordine finito). L'azione $\rho^*(g) = (\rho(g^{-1}))^T$ ha autovalori $\lambda_1^{-1}, \ldots, \lambda_n^{-1}$. Poiché $|\lambda_i| = 1$, si ha $\lambda_i^{-1} = \bar{\lambda}_i$, quindi:

$$\chi_{V^*}(g) = \sum_i \lambda_i^{-1} = \sum_i \bar{\lambda}_i = \overline{\sum_i \lambda_i} = \overline{\chi_V(g)}.$$
}

\subsection{Spazio Hom di rappresentazioni}
Usando l'isomorfismo $\mathrm{Hom}_K(V,U) \cong V^* \otimes U$, l'azione di $G$ su $\mathrm{Hom}_K(V,U)$ si ottiene combinando duale e prodotto tensoriale. Esplicitamente: per $f \in \mathrm{Hom}_K(V,U)$, definiamo $ \xi: G \to GL(Hom_K(V,U)) $ come:

$$\xi (g)(f)(v) = (\pi(g) \circ f \circ \rho(g)^{-1})(v)$$

Cioè: "coniuga $f$ con le rappresentazioni di $g$". Si verifica che questo definisce un omomorfismo.

Abbiamo ottenuto questa definizione usando l'isomorfismo canonico, per cui:
\[
  f(v) = \varphi(v)u
\]
e' l'omomorfismo corrispondente a $ f $ dati $ \varphi $ e $ u $, ai quali possiamo applicare la rappresentazione di prodotto tensoriale
\[
  (\rho^* \otimes \pi)(g)(\varphi \otimes u) = (\rho^*(g)\varphi) \otimes (\pi(g)u)
\]
e di spazio duale
\[
  (\varphi \circ \rho(g)^{-1}) \otimes (\pi(g)u)
\]
per poi trovare l'omomorfismo associato usando $ \Psi $:
\[
  \Psi((\varphi \circ \rho(g)^{-1}) \otimes (\pi(g)u))(v) = (\varphi \circ \rho(g)^{-1})(v) \cdot (\pi(g)u)
\]
dato che la parte a sx e' uno scalare e $ \pi $ e' lineare, possiamo portarlo dentro
\[
  \pi(g)(\varphi \circ \rho(g)^{-1}(v)\cdot u)
\]
dove $ \varphi(\rho(g)^{-1}(v)) \cdot u $ puo' essere riscritto come $ f(\rho(g)^{-1}(v)) $
\[
  \pi(g) \circ f \circ \rho(g)^{-1}(v)
\]
che e' proprio la definizione che abbiamo dato.

\cor{Carattere di $ Hom(V,U) $}{
  Date due rappresentazioni $ V, U $ con caratteri $ \chi_V, \chi_U $:
  \[
    \chi_{Hom(V,U)}(g) = \chi_{V^*\otimes U} = \overline{\chi_V(g)} \cdot \chi_U(g) 
  \]
}
\pf{Dimostrazione}{
  Ovvia per proposizioni precedenti.
}

\section{Ortogonalita' dei Caratteri}
\subsection{La componente isotipica banale $ V^G $}
\dfn{$ V^G $}{
Data una rappresentazione $(V, \rho)$ di $G$, la componente isotipica banale (o spazio dei punti fissi) è:

$$V^G := \{v \in V \mid \rho(g)\cdot v = v \;\;\forall g \in G\}$$  
}

È il sottospazio di $V$ su cui $G$ agisce banalmente, cioè la somma di tutte le copie della rappresentazione banale dentro $V$.

\nt{
  Le rappresentazioni \textit{banali} e \textit{triviali} sono la stessa cosa, sono semplicemente quelle rappresentazioni su cui tutti gli elementi di $ G $ agiscono come identita'.
}

  
\subsubsection{Costruzione del proiettore su $ V^G $}
Costruiamo un proiettore $ \varphi: V \to V^G $

$$\varphi := \frac{1}{|G|}\sum_{g \in G} \rho(g) \;\in\; \mathrm{End}(V)$$

è ben definita ed è un morfismo di rappresentazioni. Si verifica facilmente che:
\begin{itemize}
\item $\varphi$ è un proiettore: $\varphi^2 = \varphi$.
  \item $\mathrm{Im}(\varphi) = V^G$.
\end{itemize}
Quindi $\varphi: V \to V^G$ è la proiezione canonica su $V^G$.

\subsubsection{Dimensione di $ V^G $ e molteplicita' della rappresentazione banale}
Essendo $\varphi$ un proiettore, $\dim V^G = \dim \mathrm{Im}(\varphi)$. Inoltre, il suo rango e' uguale alla sua traccia (dato che e' diagonalizzabile e sappiamo che i suoi autovalori devono essere 0 o 1), quindi $ \dim V^G = Tr(\varphi) $. Ma:

$$\mathrm{Tr}(\varphi) = \frac{1}{|G|}\sum_{g \in G} \mathrm{Tr}(\rho(g)) = \frac{1}{|G|}\sum_{g \in G} \chi_V(g)$$

Quindi:

$$\dim V^G = \frac{1}{|G|}\sum_{g \in G} \chi_V(g)$$

Possiamo dire quindi che:
\mlenma{Lemma A}{
  La molteplicità della rappresentazione banale in $V$ è $\dim V^G = \dfrac{1}{|G|}\displaystyle\sum_{g \in G}\chi_V(g)$.
}
Quindi il numero delle sottorapresentazioni banali (di dimensione 1) di una rappresentazione e' data dalla somma dei caratteri di tutti gli elementi di $ G $ diviso per il numero di elementi in $ G $.

\subsubsection{La rappresentazione $ Hom_G(U,W) $}
Consideriamo il caso in cui $ U $ e $ W $ sono rappresentazioni di $ G $ e $ V = Hom_\mathbb{C}(U,W) $. Per quanto abbiamo visto in precedenza, anche $ V $ costituisce una rappresentazione di $ G $, quindi esiste
\[
  V^G = (Hom_\mathbb{C}(U,W))^G = \{f: V \to W \text{ lin. } | g \cdot f = f\ \forall g \in G\}
\]
dove l'azione di $g$ su è la coniugazione $g \cdot f = \pi(g) \circ f \circ \rho(g)^{-1}$ (come definito in precedenza). Ma questa e' la condizione per essere un morfismo di rappresentazioni, quindi $ (Hom_\mathbb{C}(U,W))^G $ e' anche lo spazio dei morfismi di rappresentazioni tra $ U $ e $ W $, indicato con $ Hom_G(U,W) $.

Abbiamo quindi mostrato che:

\mlenma{Lemma B}{
  \[
    (Hom_\mathbb{C}(U,W))^G = Hom_G(U,W)
  \]
}

Quindi la componente isotipica banale dello spazio degli omomorfismi da due rappresentazioni di $ G $ e' l'insieme degli omomorfismi di rappresentazioni da una rappresentazione all'altra.

\subsection{Funzioni Centrali e Prodotto Hermitiano}
\dfn{Spazio delle funzioni centrali}{
  Lo spazio delle funzioni centrali su $ G $ e':
  \[ 
\mathcal{C}(G) \coloneq \{f: G \to \mathbb{C} \mid f(hgh^{-1}) = f(g) \;\;\forall g,h \in G\}
  \]
}
 I caratteri $\chi_V$ sono funzioni centrali (per ciclicità della traccia). La dimensione di $\mathcal{C}(G)$ è uguale al numero di classi di coniugio di $G$, poiché una funzione centrale è determinata dai suoi valori su ciascuna classe.\\

 Definiamo un prodotto scalare (hermitiano dato che $ K = \mathbb{C} $) su questo spazio:
\dfn{Prodotto hermitiano fra funzioni centrali}{
  Il prodotto hermitiano su $\mathcal{C}(G)$ è:
  $$\langle \chi_1, \chi_2 \rangle := \frac{1}{|G|} \sum_{g \in G} \chi_1(g)\,\overline{\chi_2(g)}$$ 
 }
Si verifica che è lineare nel primo argomento, antilineare nel secondo (per via del coniugato), e definito positivo. È quindi un vero prodotto hermitiano su $\mathcal{C}(G)$.

\subsection{Teorema di Ortogonalita'}
\thm{Ortogonalita' dei Caratteri}{
  Siano $U, W$ rappresentazioni irriducibili di $G$. Allora:
  $$\langle \chi_U, \chi_W \rangle = \begin{cases} 1 & \text{se } U \cong W \\ 0 & \text{se } U \not\cong W \end{cases}$$ 
}
\pf{Dimostrazione}{
  Consideriamo la rappresentazione $\mathrm{Hom}_K(U,W) \cong U^* \otimes W$. Il suo carattere è $\chi_{U^* \otimes W} = \overline{\chi_U} \cdot \chi_W$. Quindi:
  $$\langle \chi_U, \chi_W \rangle = \frac{1}{|G|}\sum_{g \in G} \chi_W(g)\,\overline{\chi_U(g)} = \frac{1}{|G|}\sum_{g \in G}\chi_{U^* \otimes W}(g) \overset{\text{Lemma A}}{=} \dim(\mathrm{Hom}_K(U,W))^G$$
%
  Per il Lemma B:
    $$= \dim \mathrm{Hom}_G(U,W)$$
  %
  Ora applichiamo il Lemma di Schur:
  \begin{itemize}
    \item Se $U \not\cong W$: ogni morfismo $f: U \to W$ di rappresentazioni è zero, quindi $\mathrm{Hom}_G(U,W) = \{0\}$, dunque $\dim = 0$.
    \item Se $U \cong W$: i morfismi di rappresentazioni $U \to U$ sono solo le omotetie $\lambda \cdot \mathrm{id}$ (spazio 1-dimensionale), quindi $\dim \mathrm{Hom}_G(U,U) = 1$.
  \end{itemize}
  %
  Catena logica \textit{limpida}
  \[
    \langle \chi_U, \chi_W \rangle \overset{\text{A}}{=} \dim(\mathrm{Hom}_K(U,W))^G \overset{\text{B}}{=} \dim\mathrm{Hom}_G(U,W) \overset{\text{Schur}}{=} \begin{cases}1 \\ 0\end{cases}
  \]
}

\ex{Verifica con $S_3$}{
Prendiamo $U =$ rappresentazione banale e $W =$ rappresentazione standard di $S_3$ (dim 2). Le classi di coniugio sono $\{e\}$, $\{(12),(13),(23)\}$, $\{(123),(132)\}$ con $|C_i| = 1, 3, 2$.

  \begin{center}
    \begin{tabular}{c|cc}
      Classe & $\chi_U$ & $\chi_W$ \\
      \hline
      $e$ & $1$ & $2$ \\
      trasposizioni & $1$ & $0$ \\
      3-cicli & $1$ & $-1$
    \end{tabular}
  \end{center}

  Poiché i caratteri sono reali per $S_3$, si ha $\overline{\chi_W} = \chi_W$. Quindi:
  $$\langle \chi_U, \chi_W \rangle = \frac{1}{6}(1\cdot2\cdot1 + 1\cdot0\cdot3 + 1\cdot(-1)\cdot2) = \frac{1}{6}(2+0-2) = 0 \quad \checkmark$$
  $$\langle \chi_W, \chi_W \rangle = \frac{1}{6}(4\cdot1 + 0\cdot3 + 1\cdot2) = \frac{1}{6}(4+0+2) = 1 \quad \checkmark$$
}

\subsection{Decomposizione e Componenti Isotipiche}

\thm{Decomposizione Isotipica}{
  Siano $V_1, \ldots, V_r$ tutte le rappresentazioni irriducibili di $G$ (a meno di isomorfismo). Allora ogni rappresentazione $V$ di $G$ si decompone in modo unico come:
  $$V \cong V_1^{\oplus k_1} \oplus \cdots \oplus V_r^{\oplus k_r}$$
  dove $V_i^{\oplus k_i} = \underbrace{V_i \oplus \cdots \oplus V_i}_{k_i}$ si chiama \textbf{componente isotipica} di tipo $V_i$ e raccoglie tutte le copie isomorfe a $V_i$ in $V$.

  In termini di caratteri: $\chi_V = \sum_{i=1}^r k_i \chi_{V_i}$.
}

\cor{Calcolo delle molteplicità}{
  Le molteplicità $k_i$ sono:
  $$k_i = \langle \chi_V, \chi_{V_i} \rangle$$
}

\pf{Dimostrazione}{
  Per linearità del prodotto hermitiano e ortogonalità:
  $$\langle \chi_V, \chi_{V_j}\rangle = \left\langle \sum_i k_i \chi_{V_i}, \chi_{V_j}\right\rangle = \sum_i k_i \langle \chi_{V_i}, \chi_{V_j}\rangle = \sum_i k_i \delta_{ij} = k_j. \quad \square$$
}

\cor{Criterio di irriducibilità}{
  $V$ è irriducibile se e solo se $\langle \chi_V, \chi_V \rangle = 1$.
}

\pf{Dimostrazione}{
  $\langle \chi_V, \chi_V \rangle = \sum_i k_i^2$. Poiché $k_i \in \mathbb{Z}_{\geq 0}$, si ha $\sum_i k_i^2 = 1 \iff$ esattamente un $k_i = 1$ e tutti gli altri sono 0. $\square$
}

\cor{Il carattere è un invariante completo}{
  $V \cong W \iff \chi_V = \chi_W$.
}

\pf{Dimostrazione}{
  $(\Rightarrow)$ Ovvio per definizione di carattere. $(\Leftarrow)$ Se $\chi_V = \chi_W$, le molteplicità di ogni irriducibile in $V$ e in $W$ coincidono per il corollario precedente, quindi $V \cong W$. $\square$
}

\nt{
Due rappresentazioni sono isomorfe se e solo se hanno lo stesso carattere: il carattere è dunque un \textbf{invariante completo} della rappresentazione.
}

\subsection{Struttura della Rappresentazione Regolare}

\thm{Molteplicità nella rappresentazione regolare}{
  Ogni rappresentazione irriducibile $V_i$ appare in $\mathbb{C}[G]$ con molteplicità $k_i = \dim V_i$.
}

\pf{Dimostrazione}{
  Per il corollario sul calcolo delle molteplicità:
  $$k_i = \langle \chi_{V_i}, \chi_{\mathbb{C}[G]} \rangle = \frac{1}{|G|}\sum_{g \in G} \chi_{V_i}(g)\,\overline{\chi_{\mathbb{C}[G]}(g)}$$
  Il carattere regolare vale $\chi_{\mathbb{C}[G]}(g) = \begin{cases}|G| & g = e \\ 0 & g \neq e\end{cases}$, quindi:
  $$k_i = \frac{1}{|G|} \cdot \chi_{V_i}(e) \cdot |G| = \chi_{V_i}(e) = \dim V_i. \quad \square$$
}

\cor{Identità di Burnside}{
  Poiché $\dim \mathbb{C}[G] = |G|$:
  $$|G| = \sum_{i=1}^r (\dim V_i)^2$$
}

\subsection{I Caratteri Irriducibili come Base di $\mathcal{C}(G)$}

Per dimostrare che i caratteri irriducibili formano una \textit{base} (non solo un sistema ortonormale) di $\mathcal{C}(G)$, introduciamo uno strumento ausiliario.

\dfn{Operatore $\psi_{\alpha,V}$}{
  Sia $(V, \rho)$ una rappresentazione di $G$ e $\alpha \in \mathcal{C}(G)$ una funzione centrale. Definiamo:
  $$\psi_{\alpha, V} := \sum_{g \in G} \alpha(g)\, \rho(g) \;\;\in\;\; \mathrm{End}_K(V)$$
  È un operatore su $V$ ottenuto pesando le trasformazioni $\rho(g)$ con i valori di $\alpha$.
}

\mlenma{$\psi_{\alpha,V}$ è un morfismo di rappresentazioni}{
  $\psi_{\alpha,V} \in \mathrm{End}_G(V)$, ovvero commuta con l'azione di $G$.
}

\pf{Dimostrazione}{
  Sia $h \in G$. Calcoliamo $\psi_{\alpha,V} \circ \rho(h)$: sostituendo $s = gh$ (biiezione di $G$),
  $$\psi_{\alpha,V}\circ\rho(h) = \sum_{g \in G} \alpha(g)\,\rho(gh) = \sum_{s \in G} \alpha(sh^{-1})\,\rho(s).$$
  Calcoliamo $\rho(h) \circ \psi_{\alpha,V}$: sostituendo $s = hg$,
  $$\rho(h) \circ \psi_{\alpha,V} = \sum_{g \in G} \alpha(g)\,\rho(hg) = \sum_{s \in G} \alpha(h^{-1}s)\,\rho(s).$$
  Poiché $\alpha$ è centrale: $\alpha(sh^{-1}) = \alpha(h^{-1}s)$ (ciclicità). Quindi i due lati coincidono. $\square$
}

\nt{
  Se $V$ è irriducibile, per il Lemma di Schur $\psi_{\alpha,V} = \lambda_{\alpha,V} \cdot \mathrm{id}_V$. Prendendo la traccia:
  $$\lambda_{\alpha,V} \cdot \dim V = \mathrm{tr}(\psi_{\alpha,V}) = \sum_{g \in G}\alpha(g)\chi_V(g)$$
  quindi:
  $$\lambda_{\alpha,V} = \frac{1}{\dim V}\sum_{g \in G}\alpha(g)\chi_V(g)$$
}

\thm{I caratteri irriducibili sono una base di $\mathcal{C}(G)$}{
  I caratteri $\chi_{V_1}, \ldots, \chi_{V_r}$ delle rappresentazioni irriducibili (a meno di isomorfismo) formano una \textbf{base ortonormale} di $\mathcal{C}(G)$.
}

\pf{Dimostrazione}{
  \textbf{Indipendenza lineare}: segue dall'ortogonalità — formano già un sistema ortonormale.

  \textbf{Generazione}: mostriamo che se $\alpha \in \mathcal{C}(G)$ soddisfa $\langle \alpha, \chi_{V_i}\rangle = 0$ per ogni $i$, allora $\alpha = 0$.

  Per ipotesi, $\lambda_{\alpha,V_i} = \frac{1}{\dim V_i}\sum_g \alpha(g)\chi_{V_i}(g) = 0$, quindi $\psi_{\alpha,V_i} = 0$ per ogni irriducibile $V_i$.

  Per additività ($\psi_{\alpha,V\oplus W} = \psi_{\alpha,V}\oplus\psi_{\alpha,W}$), si ha $\psi_{\alpha,V} = 0$ per ogni rappresentazione $V$.

  In particolare, per la rappresentazione regolare $\mathbb{C}[G]$:
  $$\psi_{\alpha,\mathbb{C}[G]}(e_e) = \sum_{g \in G}\alpha(g)\,\rho_{\mathrm{reg}}(g)(e_e) = \sum_{g \in G}\alpha(g)\, e_g = 0$$
  Ma $\{e_g\}_{g\in G}$ è una base di $\mathbb{C}[G]$, quindi tutti i coefficienti $\alpha(g) = 0$, ovvero $\alpha = 0$. $\square$
}

\section{Corollari Strutturali}

\cor{Numero di rappresentazioni irriducibili}{
  Il numero di rappresentazioni irriducibili di $G$ (a meno di isomorfismo) è uguale al numero di classi di coniugio di $G$.
}

\pf{Dimostrazione}{
  I caratteri irriducibili formano una base di $\mathcal{C}(G)$, che ha dimensione uguale al numero di classi di coniugio. $\square$
}

\cor{Rappresentazioni irriducibili di gruppi abeliani}{
  $G$ è abeliano $\iff$ tutte le sue rappresentazioni irriducibili hanno dimensione 1.
}

\pf{Dimostrazione}{
  \begin{itemize}
    \item $(\Rightarrow)$: Se $G$ è abeliano, ogni elemento ha la propria classe di coniugio, quindi ci sono $|G|$ classi e di conseguenza $|G|$ irriducibili. Per Burnside: $\sum_i d_i^2 = |G|$ con $|G|$ addendi implica $d_i = 1$ per ogni $i$.

    \item $(\Leftarrow)$: Se tutti i $d_i = 1$, dalla decomposizione $\mathbb{C}[G] \cong \mathbb{C}^{|G|}$ si deduce che $\mathbb{C}[G]$ è commutativo, da cui $G$ è abeliano. $\square$
  \end{itemize}
}

\section{Esempio Completo: $S_3$}

$S_3$ ha 3 classi di coniugio, quindi esattamente 3 rappresentazioni irriducibili $V_1, V_2, V_3$.

Da Burnside: $d_1^2 + d_2^2 + d_3^2 = |S_3| = 6$. Con $d_1 = 1$ (banale $U$) e $d_2 = 2$ (standard $W$), si ottiene $d_3^2 = 6 - 1 - 4 = 1$, dunque $d_3 = 1$ (rappresentazione segno $X$).

\begin{center}
\begin{tabular}{c|c|ccc}
Classe & $|C_i|$ & $\chi_U$ & $\chi_W$ & $\chi_X$ \\
\hline
$e$ & 1 & 1 & 2 & 1 \\
trasposizioni & 3 & 1 & 0 & $-1$ \\
3-cicli & 2 & 1 & $-1$ & 1
\end{tabular}
\end{center}

La decomposizione della rappresentazione regolare è:
$$\mathbb{C}[S_3] \cong U^{\oplus 1} \oplus W^{\oplus 2} \oplus X^{\oplus 1}$$
con $1^2 + 2^2 + 1^2 = 1 + 4 + 1 = 6 = |S_3|$


%
% \end{document}
